\section{The Interaction Picture and Rabi oscillation (11/20)}

\subsection{The interaction picture}
Let
\begin{equation}
 \hat H(t)=\hat H_0+\hat V(t),
 \qquad
 \hat H_0\dket a=E_a\dket a.
\end{equation}
The Schrödinger-picture evolution operator satisfies
\begin{equation}
 \imth\hbar\frac{\partial}{\partial t}\hat U(t,t_0)
 =\hat H(t)\hat U(t,t_0).
\end{equation}
Define
\begin{align}
 \dket{\psi_I(t)}
 &=e^{\imth\hat H_0t/\hbar}\dket{\psi_S(t)},\\
 \hat U_I(t,t_0)
 &=e^{\imth\hat H_0t/\hbar}
 \hat U(t,t_0)e^{-\imth\hat H_0t_0/\hbar},\\
 \hat V_I(t)
 &=e^{\imth\hat H_0t/\hbar}
 \hat V(t)e^{-\imth\hat H_0t/\hbar}.
\end{align}
Then
\begin{align}
 \imth\hbar\frac{\dif}{\dif t}\dket{\psi_I(t)}
 &=\hat V_I(t)\dket{\psi_I(t)},\\
 \imth\hbar\frac{\partial}{\partial t}\hat U_I(t,t_0)
 &=\hat V_I(t)\hat U_I(t,t_0).
\end{align}
An interaction-picture operator is
\begin{equation}
 \hat O_I(t)=e^{\imth\hat H_0t/\hbar}
 \hat O_Se^{-\imth\hat H_0t/\hbar},
\end{equation}
with
\begin{equation}
 \imth\hbar\frac{\dif\hat O_I}{\dif t}
 =\comm{\hat O_I}{\hat H_0}.
\end{equation}

\subsection{Dyson equation and Dyson series}
Integrating the evolution equation gives
\begin{equation}
 \boxed{
 \hat U_I(t,t_0)
 =\id-\frac{\imth}{\hbar}
 \int_{t_0}^{t}\hat V_I(t')\hat U_I(t',t_0)\dif t'.}
\end{equation}
Iteration yields
\begin{align}
 \hat U_I(t,t_0)
 &=\id-\frac{\imth}{\hbar}\int_{t_0}^{t}
 \hat V_I(t_1)\dif t_1+
 \left(-\frac{\imth}{\hbar}\right)^2
 \int_{t_0}^{t}\dif t_1
 \int_{t_0}^{t_1}\dif t_2\,
 \hat V_I(t_1)\hat V_I(t_2)+\cdots\\
 &=\mathcal T\exp\left[-\frac{\imth}{\hbar}
 \int_{t_0}^{t}\hat V_I(t')\dif t'\right].
\end{align}

\subsubsection{Adiabatic switching and Gellmann-Low theorem}
We start from the following Hamiltonian, where the interaction $\hat V$ is adiabatically switched on in the $\epsilon\rightarrow0^+$ limit. 
\begin{equation}
 \hat H(t)=\hat H_0+\hat V e^{-\epsilon|t|},
 \qquad -\infty<t<0,
 \qquad \epsilon>0.
\end{equation}
\begin{thm} \textbf{Gellmann-Low theorem:} If $\dket a$ is an eigenstate of $\hat H_0$, then
\begin{equation}
 \boxed{
 \dket{\psi_a^+}
 =\lim_{\epsilon\to0^+}
 \lim_{t_0\to-\infty}
 \hat U_I(0,t_0)\dket a}
\end{equation}
is an eigenstate of $\hat H_0+\hat V$, up to normalization and the usual
adiabatic phase.
\end{thm}
 This is known as the Gellmann-Low theorem. It can be derived from Dyson's equation and the Lippmann--Schwinger equation.

\subsection{Rabi Oscillations}

\subsubsection{Spin one-half in a time-dependent magnetic field}
Let
\begin{equation}
 \hat H(t)=-\bm\mu\cdot\bm B(t),
 \qquad
 \bm\mu=\frac{ge}{2m_e}\bm S,
\end{equation}
and
\begin{equation}
 \bm B(t)=B_0\hat{\bm z}
 +B_1\left(\hat{\bm x}\cos\omega t
 -\hat{\bm y}\sin\omega t\right).
\end{equation}
Write
\begin{equation}
 \hat H(t)=\hat H_0+\hat V(t),
 \qquad
 \hat H_0=-\frac{\hbar\omega_{12}}{2}\sigma_z,
\end{equation}
where
\begin{equation}
 \omega_{12}=\frac{geB_0}{2m_e},~~~\gamma=\frac{geB_1}{4m_e}
\end{equation}
Let
\begin{equation}
 \hat V(t)=\hbar\gamma
 \left(\sigma_x\cos\omega t-\sigma_y\sin\omega t\right)
 =\hbar\gamma
 \begin{pmatrix}
 0&e^{\imth\omega t}\\
 e^{-\imth\omega t}&0
 \end{pmatrix}.
\end{equation}
In the interaction picture,
\begin{equation}
 \hat V_I(t)=-\hbar\gamma
 \begin{pmatrix}
 0&e^{\imth(\omega-\omega_{12})t}\\
 e^{-\imth(\omega-\omega_{12})t}&0
 \end{pmatrix}.
\end{equation}

\subsection{Coupled amplitudes}
Write
\begin{equation}
 \dket{\psi_I(t)}
 =C_1(t)\dket\uparrow+C_2(t)\dket\downarrow.
\end{equation}
Then
\begin{align}
 \imth\dot C_1
 &=-\gamma e^{\imth(\omega-\omega_{12})t}C_2,\\
 \imth\dot C_2
 &=-\gamma e^{-\imth(\omega-\omega_{12})t}C_1.
\end{align}
With detuning
\begin{equation}
 \Delta=\omega-\omega_{12},
\end{equation}
one obtains
\begin{equation}
 \ddot C_1-\imth\Delta\dot C_1+\gamma^2C_1=0.
\end{equation}
The characteristic frequencies are
\begin{equation}
 \Omega_\pm
 =\frac{\Delta}{2}
 \pm\sqrt{\gamma^2+\frac{\Delta^2}{4}}.
\end{equation}
Define the generalized Rabi frequency
\begin{equation}
 \boxed{\Omega_R=\sqrt{\gamma^2+\frac{\Delta^2}{4}}.}
\end{equation}

For
\begin{equation}
 C_1(0)=1,
 \qquad C_2(0)=0,
\end{equation}
the transition probability is
\begin{equation}
 \boxed{
 |C_2(t)|^2
 =\frac{\gamma^2}{\gamma^2+\Delta^2/4}
 \sin^2\left(\Omega_Rt\right).}
\end{equation}
The survival probability is
\begin{equation}
 |C_1(t)|^2=1-|C_2(t)|^2.
\end{equation}
At resonance,
\begin{equation}
 \omega=\omega_{12},~~~\Delta=0
\end{equation}
we have 
\begin{equation}
 \boxed{|C_2(t)|^2=\sin^2(\gamma t),
 \qquad |C_1(t)|^2=\cos^2(\gamma t).}
\end{equation}

\subsection{Nuclear magnetic resonance}
The resonance condition is
\begin{equation}
 \boxed{\omega_{\mathrm{res}}=\omega_{12}=\mu B_0.}
\end{equation}
In an NMR experiment,
\begin{equation}
 B_0=H+M=H+\chi H.
\end{equation}
The magnetization and susceptibility are
\begin{equation}
 M=\frac{\omega_{\mathrm{res}}}{\mu}-H,
 \qquad
 \chi=\frac{M}{H}=\frac{\omega_\text{res}}{\mu H}-1.
\end{equation}